\documentclass[11pt]{article}

\usepackage{series}
\usepackage{amsmath,amssymb,amsfonts,mathtools}
\usepackage{hyperref}
\usepackage{enumitem}

\newcommand{\Hilb}{\mathcal H}
\newcommand{\Dens}{\mathsf D}
\newcommand{\Bas}{\mathcal B}
\newcommand{\Inst}{\mathcal I}
\newcommand{\Chan}{\mathcal D}
\newcommand{\Tr}{\operatorname{Tr}}

\title{\bf Why Decoherence Cannot Replace Measurement\\
\large Outcome-Resolved Dynamics Without a Fundamental Measurement Postulate
in Modal Triplet Theory}
\author{Peter Nero}
\date{July 2026}

\begin{document}
\maketitle

\begin{abstract}
Measurement is not taken to be a fundamental act, an observer intervention,
or a separate kind of physics.  It is an ordinary
system--apparatus--environment process described with a retained physical
record.  If that record is discarded, the same coupling is represented only
by a nonselective channel.  Decoherence can suppress interference and stabilize
preferred records, but the record-discarded channel does not contain an
outcome label or a conditional state update.  We make this data distinction
exact in finite-dimensional quantum theory and then state its consequence for
Modal Triplet Theory (MTT).  For a pointer resolution
$\{P_a\}$, the dephasing channel
$\Delta_P(\rho)=\sum_aP_a\rho P_a$ is completely positive and trace
preserving, removes every off-diagonal pointer block, and preserves all pointer
populations.  We prove that the same nonselective channel is compatible with
distinct outcome instruments having different probability laws.  Hence the
channel alone specifies neither records nor their probabilities.

An MTT basin boundary or chart exit has the same logical limitation: it marks
the failure of a supplied effective chart, but it does not choose a successor
basin.  A complete outcome-resolved model must additionally provide an
outcome-indexed completely positive instrument or a normalized
selection-completion kernel, together with its source and domain.  This yields
a rigorous conditional bridge between record-discarded intra-basin suppression
and record-resolved inter-basin continuation.  Both are ordinary physical
dynamics; the distinction concerns which data the effective description
retains.  The
current canonical q79 binary recorder supplies an exact stopped-output law on
its declared finite domain, but general apparatus contexts and objective
single-history actualization remain open.  The paper therefore retains the
decoherence-versus-selection distinction while withdrawing any inference of
capture, Born probabilities, irreversibility, or objective outcomes from
projection, noninjectivity, admissibility loss, or chart exit alone.
\end{abstract}

\tableofcontents

\section{Scope and Claim Tier}

Throughout this paper, ``measurement'' names no privileged physical process.
It means an ordinary interaction whose macroscopic record is included in the
description.  The title's phrase ``cannot replace'' is therefore a statement
about information: a record-discarded marginal cannot provide the
outcome-resolved data that were discarded.

Environment-induced decoherence is a dynamical account of interference
suppression and pointer stability~\cite{Zurek2003,Schlosshauer2005}.  The
definite-outcome problem asks for something else: why an individual run has a
recorded outcome, what its conditional post-measurement state is, and what law
governs the alternatives.  Conflating these questions obscures both the
success of decoherence and the data still required for measurement.

This revision separates four claim tiers.
\begin{enumerate}[label=\arabic*.]
\item \textbf{Established measurement mathematics.}  Channels describe
nonselective dynamics; instruments describe outcome probabilities and
conditional updates~\cite{DaviesLewis1970}.
\item \textbf{Exact result of this paper.}  A dephasing channel does not
determine an instrument, and loss of an effective chart does not determine a
continuation kernel.
\item \textbf{Conditional MTT bridge.}  If a selected MTT descent emits
basins, basin-preserving channels, and an inter-basin completion instrument,
then record-discarded and record-resolved descriptions can be represented as
intra- and inter-basin operations of the same physical system.
\item \textbf{Current selected finite-domain result.}  The canonical q79
binary one-anchor recorder emits an operational stopped-output measure and
second-moment capture descent on its declared domain
\cite{NeroQ79Recorder}.  This is not yet a universal apparatus theorem or an
objective single-history theorem.
\end{enumerate}

No result below derives a Born source, physical time, entropy arrow, or
objective actualization from projection alone.  The analysis is
finite-dimensional so that every map and norm is elementary.  The distinction
extends to normal instruments on operator algebras, but that extension is not
needed here.

\section{Channels, Pointer Blocks, and Decoherence}

Let $\Hilb$ be a finite-dimensional Hilbert space and let
\[
 \Dens(\Hilb)=\{\rho\geq0:\Tr\rho=1\}.
\]
A finite pointer resolution is a family of mutually orthogonal projections
$P=\{P_a\}_{a\in A}$ satisfying $\sum_aP_a=I$.

\begin{definition}[Nonselective dephasing channel]
The $P$-dephasing channel is
\begin{equation}
 \Delta_P(\rho)=\sum_{a\in A}P_a\rho P_a.
 \label{eq:dephasing}
\end{equation}
Its pointer-coherence functional is
\begin{equation}
 C_P(\rho)^2
 =\sum_{\substack{a,b\in A\\a\neq b}}
   \lVert P_a\rho P_b\rVert_2^2,
 \label{eq:coherence}
\end{equation}
where $\lVert X\rVert_2^2=\Tr(X^\dagger X)$.
\end{definition}

\begin{proposition}[What dephasing does]\label{prop:dephasing}
The map $\Delta_P$ is completely positive and trace preserving.  For every
$\rho\in\Dens(\Hilb)$,
\[
 C_P(\Delta_P(\rho))=0,
 \qquad
 \Tr(P_a\Delta_P(\rho))=\Tr(P_a\rho)
\]
for all $a\in A$.
\end{proposition}

\begin{proof}
Equation~\eqref{eq:dephasing} is a Kraus representation with Kraus operators
$P_a$, and $\sum_aP_a^\dagger P_a=I$.  Orthogonality gives
$P_a\Delta_P(\rho)P_b=0$ for $a\neq b$, proving the first identity.  Cyclicity
of the trace and $P_aP_b=\delta_{ab}P_a$ give the second.
\end{proof}

More general decoherence dynamics may approach $\Delta_P$ only
asymptotically, may include dissipation, and may have an approximate pointer
resolution.  None of those refinements changes the type of the object:
discarding the environmental record produces a nonselective quantum channel.
It produces a density operator, not a classical outcome label.

\begin{example}[Suppression without an outcome]\label{ex:qubit}
For $\Hilb=\mathbb C^2$, let $P_0=\lvert0\rangle\langle0\rvert$,
$P_1=\lvert1\rangle\langle1\rvert$, and
$\lvert+\rangle=(\lvert0\rangle+\lvert1\rangle)/\sqrt2$.  Then
\[
 \Delta_P(\lvert+\rangle\langle+\rvert)
 =\frac12P_0+\frac12P_1.
\]
The off-diagonal blocks vanish, but the output is not either conditional
state $P_0$ or $P_1$, and no recorded value $0$ or $1$ appears in the channel
output.
\end{example}

\section{Outcome-Resolved Processes Require an Instrument}

An ordinary physical coupling can be represented at two resolutions.  If its
classical record is discarded, it gives a nonselective channel.  If the record
is retained, it gives an outcome-resolved instrument.

\begin{definition}[Quantum instrument]\label{def:instrument}
A finite quantum instrument is a family
$\{\Inst_a\}_{a\in A}$ of completely positive, trace-nonincreasing maps such
that
\[
 \Inst_{\mathrm{ns}}=\sum_{a\in A}\Inst_a
\]
is trace preserving.  For input $\rho$ it defines
\begin{equation}
 p(a\mid\rho)=\Tr\Inst_a(\rho),
 \qquad
 \rho_a=\frac{\Inst_a(\rho)}{p(a\mid\rho)}
 \quad\text{when }p(a\mid\rho)>0.
 \label{eq:instrument-law}
\end{equation}
\end{definition}

The instrument contains three pieces absent from the record-discarded channel:
an outcome index, its probability, and its conditional update.  This does not
introduce a fundamental measurement postulate.  It is the outcome-resolved
bookkeeping of the same kind of physical interaction.  If one additionally
claims that exactly one ontic history is actualized, that stronger ontology
must identify what constitutes the realized event; it is not required for the
ordinary operational meaning of a detector record.

\begin{theorem}[A nonselective channel does not select an
instrument]\label{thm:underdetermination}
The dephasing channel $\Delta_P$ does not determine outcome probabilities or
conditional states.  In particular, for the two-outcome qubit resolution of
Example~\ref{ex:qubit}, the L\"uders instrument
\[
 \Inst_0(\rho)=P_0\rho P_0,\qquad
 \Inst_1(\rho)=P_1\rho P_1
\]
and, for every $q\in[0,1]$, the instrument
\[
 \mathcal J^{(q)}_0(\rho)=q\Delta_P(\rho),\qquad
 \mathcal J^{(q)}_1(\rho)=(1-q)\Delta_P(\rho)
\]
have the same nonselective channel $\Delta_P$.  On
$\lvert+\rangle\langle+\rvert$ their outcome laws are respectively
$(1/2,1/2)$ and $(q,1-q)$.
\end{theorem}

\begin{proof}
Each displayed map is completely positive and trace nonincreasing.  Both
families sum to $\Delta_P$.  Their probabilities follow from
Equation~\eqref{eq:instrument-law}.  Choosing $q\neq1/2$ proves that the same
nonselective channel is compatible with different outcome laws.  Their
conditional states also differ: the L\"uders instrument gives $P_a$, whereas
$\mathcal J^{(q)}_a$ gives the same dephased state for either nonzero outcome.
\end{proof}

\begin{corollary}[Decoherence non-selection]\label{cor:nonselection}
No outcome law, conditional update, or unique realized outcome can be inferred
from a decoherence channel without additional outcome-resolved data.
\end{corollary}

The artificial family $\mathcal J^{(q)}$ is not proposed as a physical
detector.  Its purpose is logical: after the record is discarded, the channel
does not encode which outcome-resolved process produced it.  A source theorem
must select the detector coupling, readout algebra, and instrument
independently of the observed target probabilities.

\section{The Missing Selection-Completion Map}

The basin language used in MTT can be made precise without turning chart
failure into an outcome theorem.  Let $U\subseteq\Dens(\Hilb)$ be the domain
of an effective chart and let $\{\Bas_a\}_{a\in A}$ be declared measurable
successor regions.

\begin{definition}[Selection-completion kernel]\label{def:kernel}
For a boundary state $\eta\in\partial U$, a selection-completion kernel is a
probability kernel
\[
 K_\partial(a,d\sigma\mid\eta)
\]
on $A\times\Dens(\Hilb)$ satisfying
\begin{align}
 &\sum_{a\in A}\int_{\Dens(\Hilb)}
 K_\partial(a,d\sigma\mid\eta)=1,
 \label{eq:kernel-normalization}\\
 &K_\partial\bigl(a,\Dens(\Hilb)\setminus\Bas_a\mid\eta\bigr)=0.
 \label{eq:kernel-support}
\end{align}
Its nonselective continuation is the barycenter
\begin{equation}
 T_K(\eta)=
 \sum_{a\in A}\int_{\Dens(\Hilb)}
 \sigma\,K_\partial(a,d\sigma\mid\eta).
 \label{eq:barycenter}
\end{equation}
\end{definition}

A deterministic continuation is the special case in which the kernel is a
point mass.  An affine completely positive realization is more strongly
encoded by an instrument.  Definition~\ref{def:kernel} is deliberately
minimal: it completes the effective mathematical description; it does not add
a special physical act.  It exposes exactly the data that the phrase
``capture into a basin'' must supply.

\begin{proposition}[Chart exit does not determine continuation]\label{prop:chart}
The statement that a trajectory reaches $\partial U$ does not determine
$K_\partial$.  If two successor regions $\Bas_0,\Bas_1$ contain states
$\sigma_0,\sigma_1$, then both
\[
 K^{(0)}_\partial(a,d\sigma\mid\eta)
 =\delta_{a0}\delta_{\sigma_0}(d\sigma)
\]
and
\[
 K^{(1)}_\partial(a,d\sigma\mid\eta)
 =\delta_{a1}\delta_{\sigma_1}(d\sigma)
\]
satisfy Equations~\eqref{eq:kernel-normalization} and
\eqref{eq:kernel-support}, while producing different outcomes and
continuations.
\end{proposition}

\begin{proof}
Both kernels are normalized point masses supported in their declared
successor regions.  Their outcome marginals and barycenters are different.
Boundary membership supplies no condition that distinguishes them.
\end{proof}

Thus admissibility loss has a precise, limited meaning: the current effective
chart no longer controls the continuation.  It does not imply that projection
``enforces'' one particular capture.  A selected action, detector coupling,
boundary condition, instrument, or kernel must do that work.

\section{A Correct Conditional MTT Bridge}

The useful core of the original paper can now be stated without overclaiming.

\begin{definition}[Typed basin model]\label{def:basin}
A typed basin model consists of:
\begin{enumerate}[label=(\roman*)]
\item a state space and declared basin regions $\Bas_a$;
\item within-basin channels $\Chan_a$ satisfying
  $\Chan_a(\Bas_a)\subseteq\Bas_a$ on their stated domains;
\item a pointer or record algebra;
\item a selection-completion instrument or kernel at every claimed
  inter-basin event; and
\item a source certificate identifying how these objects descend from the
  same selected MTT data.
\end{enumerate}
\end{definition}

\begin{theorem}[Conditional record-discarded--record-resolved bridge]
\label{thm:bridge}
In a typed basin model, a channel $\Chan_a$ that suppresses off-diagonal
pointer blocks while preserving $\Bas_a$ is an intra-basin operation.  A
completion kernel with support in $\Bas_b$, $b\neq a$, is an inter-basin
operation.  These may be two resolutions of one ordinary physical
system--apparatus process.  The record-discarded channel does not determine the
record-resolved kernel.
\end{theorem}

\begin{proof}
The first two statements are immediate from the declared domains and support
conditions.  The last follows from
Theorem~\ref{thm:underdetermination} at channel level and
Proposition~\ref{prop:chart} at chart-boundary level.
\end{proof}

This theorem retains the intended distinction:
\[
\boxed{
\begin{aligned}
\text{record discarded} &\longleftrightarrow
  \text{intra-basin suppression and record stabilization},\\
\text{record retained} &\longleftrightarrow
  \text{an outcome-indexed inter-basin instrument or kernel}.
\end{aligned}}
\]
The arrow is a typed representation, not an ontological division into
``measurement physics'' and ordinary physics, and not a proof that both sides
have already been emitted by one universal MTT source.

\section{What Current MTT Results Supply}

Several earlier formulations attributed too much to noninvertible projection.
A noninjective map can identify microscopic states without selecting a time
arrow, an entropy law, a probability measure, or a representative of each
fiber.  The corrected MTT shadow-bridge analysis states exact descent and
restricted-recovery criteria and treats measurement instruments as selected
updates, not inverse maps~\cite{NeroProjectionProbability}.

There is nevertheless a concrete positive result.  On the canonical q79
binary one-anchor Fock recorder, the selected normal state on the commuting
nondemolition output algebra induces a stopped output law, and
second-moment capture descent is exact on the declared finite-symbol
domain~\cite{NeroQ79Recorder}.  No observed probability or fitted stochastic
parameter is added on that domain.

That theorem closes an operational finite-domain row; it does not prove:
\begin{itemize}
\item the corresponding construction for every allowed apparatus context;
\item finite-bandwidth and non-Markov corrections;
\item a pre-quantum probability semantics, if one is demanded;
\item objective actualization of one ontic history; or
\item that every admissibility boundary carries the same completion law.
\end{itemize}
Those statements require their own same-source theorems.

\section{Relation to Standard Outcome-Resolved Dynamics}

\subsection{Decoherence and pointer stability}

Standard decoherence theory derives suppression of interference in reduced
states through system--environment entanglement and explains the dynamical
selection and stability of preferred pointer structures
\cite{Zurek2003,Schlosshauer2005}.  Proposition~\ref{prop:dephasing} is the
finite idealization of this nonselective effect.  It agrees with, rather than
replaces, the standard account.

\subsection{Instruments and trajectories}

The Davies--Lewis instrument formalism associates classical outcome
statistics and conditional quantum operations with a measurement
\cite{DaviesLewis1970}.  Quantum-trajectory descriptions add a monitored
record and an unravelling of a nonselective master equation.  Different
unravellings can represent the same ensemble channel, so the master equation
alone does not select a unique record process.  This is not evidence for a
special measurement force.  It says only that MTT must emit the ordinary
detector coupling, physical readout, and outcome-resolved instrument, not
merely reproduce their record-discarded reduced channel.

\subsection{Redundant records}

Environmental redundancy can explain why information in a pointer observable
is robustly accessible.  Redundancy is compatible with the distinction made
here: it concerns the proliferation of record information, whereas
Theorem~\ref{thm:underdetermination} concerns the missing outcome-resolved law
when only the nonselective channel is supplied.  No claim that Darwinism
presupposes a particular MTT selection mechanism is needed.

\subsection{Collapse models}

Objective-collapse models append nonlinear or stochastic state dynamics that
does more than suppress ensemble coherences.  Whether an MTT limit reproduces
one such model is a separate source-and-error problem.  Even an exact
dephasing semigroup can possess a random-unitary unravelling and therefore
does not by itself prove objective collapse
\cite{NeroGravityCollapse}.

\section{Consequences and Non-Consequences}

\subsection{Born weights}

If the L\"uders instrument is supplied, then
\[
 p(a\mid\rho)=\Tr(P_a\rho)
\]
follows immediately.  This is an evaluation of a selected instrument, not a
derivation of that instrument from decoherence.  Basin sizes can represent
arbitrary finite probability vectors, so matching target weights by basin
volume is not predictive unless the measure and partition were independently
selected~\cite{NeroProjectionProbability}.

\subsection{Irreversibility}

Reduced decoherence may be effectively irreversible because environmental
records are inaccessible, and an instrument may produce persistent records.
Neither fact follows from noninjectivity alone.  An arrow of time requires a
declared state, boundary condition, semigroup, record algebra, or other
asymmetric source.

\subsection{Undecidability}

No algorithmic undecidability theorem follows from the existence of basin
boundaries.  Such a theorem requires a computably specified MTT decision
problem and an explicit reduction from a known undecidable problem.  The
earlier claim that quantum probability arises because selection is
undecidable is therefore withdrawn.

\subsection{Single outcomes}

An operational instrument predicts the distribution and conditional state of
records.  It does not by itself settle whether one ontic history is
fundamental, emergent, or interpretation-relative.  Any stronger
single-history claim must identify an actualization object and prove its
relation to the instrument.

\section{Outcome-Resolved Dynamics Promotion Contract}

A universal MTT detector-and-record theorem must emit, from one selected
source:
\begin{enumerate}[label=\arabic*.]
\item the physical state and observable algebras;
\item the pointer or record algebra and apparatus coupling;
\item the nonselective channel with its approximation domain and errors;
\item the outcome instrument or completion kernel;
\item normalization, positivity, and coarse-graining consistency;
\item the probability source without observed target weights as inputs;
\item stable record formation;
\item any claimed objective actualization rule; and
\item transport to general apparatus contexts, including finite-bandwidth and
  non-Markov regimes.
\end{enumerate}

The canonical q79 recorder supplies a nontrivial subset of this contract on
one finite domain.  The remaining rows should be attacked as source and
transport theorems, not filled by interpreting chart exit as selection.

\section{Version Delta}

Relative to version 1, this revision:
\begin{itemize}
\item retains the distinction between intra-basin suppression and inter-basin
selection;
\item states explicitly that measurement is ordinary record-forming physics,
not a fundamental act or observer-dependent postulate;
\item defines decoherence as a nonselective channel and proves its exact
suppression properties;
\item adds an explicit quantum instrument and a selection-completion kernel;
\item proves that one nonselective channel supports different instruments and
probability laws;
\item proves that chart exit does not select a successor basin or kernel;
\item replaces asserted projection-enforced capture by a conditional typed
basin theorem;
\item records the exact finite-domain q79 recorder result at its current tier;
\item withdraws Born, irreversibility, undecidability, and objective-outcome
inferences from projection or admissibility loss alone; and
\item adds a testable promotion contract for future universal measurement
claims.
\end{itemize}

\section{Conclusion}

Measurement is ordinary physics.  Nothing in this paper gives it a
fundamental status, a privileged observer, or a separate force.  The title
means only that a record-discarded description cannot supply the
outcome-resolved data it omits.  A nonselective channel can erase
off-diagonal pointer blocks and stabilize records while containing no outcome
index.  Retaining the record requires an instrument or completion kernel that
supplies outcome probabilities and conditional updates.  The same channel can
be associated with different outcome-resolved instruments, and the same chart
exit can be continued by different kernels.

This does not invalidate the MTT basin picture.  It makes the picture precise.
Intra-basin suppression and inter-basin record formation may be stages or
resolutions of one physical process once both are emitted from selected source
data.  MTT already has an exact operational account on one canonical finite
q79 recorder domain.  Extending that result to general apparatus contexts and
controlled memory effects is the operational frontier.  Objective
single-history actualization is a separate, stronger question, not part of
what makes an ordinary interaction a measurement.

\begin{thebibliography}{99}

\bibitem{DaviesLewis1970}
E.~B. Davies and J.~T. Lewis,
\emph{An operational approach to quantum probability},
Communications in Mathematical Physics \textbf{17}, 239--260 (1970).
\url{https://doi.org/10.1007/BF01647093}

\bibitem{Zurek2003}
W.~H. Zurek,
\emph{Decoherence, einselection, and the quantum origins of the classical},
Reviews of Modern Physics \textbf{75}, 715--775 (2003).
\url{https://doi.org/10.1103/RevModPhys.75.715}

\bibitem{Schlosshauer2005}
M.~Schlosshauer,
\emph{Decoherence, the measurement problem, and interpretations of quantum
mechanics},
Reviews of Modern Physics \textbf{76}, 1267--1305 (2005).
\url{https://doi.org/10.1103/RevModPhys.76.1267}

\bibitem{NeroProjectionProbability}
P.~Nero,
\emph{Projection, Probability, and Irreversibility: Shadow Bridges Between
Measurement, Black Holes, and Cosmology in Modal Triplet Theory},
corrected version 3 manuscript (2026); latest released concept record:
\url{https://doi.org/10.5281/zenodo.18256408}

\bibitem{NeroGravityCollapse}
P.~Nero,
\emph{Gravitationally Induced Collapse as an Effective Limit of Modal Triplet
Theory},
corrected version 3 manuscript (2026).

\bibitem{NeroQ79Recorder}
P.~Nero,
\emph{Canonical q79 Fock Output Measure and Second-Moment Capture Descent
Theorem},
MTT QM Source Proof repository (2026).
\url{https://github.com/PeterNero/mtt-qm-source-proof}

\bibitem{NeroFoundation}
P.~Nero,
\emph{Modal Triplet Theory: Foundation},
Zenodo (2025).
\url{https://doi.org/10.5281/zenodo.16949762}

\end{thebibliography}

\end{document}
